% Mathématiques Terminale spécialité — V3 LaTeX
% Préparation au bac — spécialité mathématiques

\section{Objectifs}
\begin{itemize}
\item Connaître la structure actuelle de l’épreuve écrite de spécialité.
\item Organiser quatre heures de travail et prioriser les questions.
\item Rédiger des raisonnements lisibles, justifiés et contrôlables.
\item Mobiliser plusieurs chapitres dans des exercices transversaux de type bac.
\end{itemize}

\section{Repères officiels de l’épreuve}
L’épreuve écrite de spécialité mathématiques dure quatre heures et l’enseignement de spécialité est affecté d’un coefficient $16$. Les sujets officiels de la session $2026$ sont constitués de quatre exercices et demandent de traiter l’ensemble du sujet. La calculatrice avec mode examen actif est autorisée, de même qu’une calculatrice sans mémoire de type collège lorsque le sujet le prévoit.

La rédaction, la clarté et la précision des raisonnements sont explicitement prises en compte. Les consignes officielles invitent à faire figurer les traces de recherche, même incomplètes ou infructueuses. Cela signifie qu’une copie ne doit pas masquer la démarche : une idée pertinente, un essai structuré ou une propriété correctement mobilisée peut contribuer à l’évaluation même si la question n’est pas entièrement achevée.

\section{1. Lecture stratégique du sujet}
Les quatre exercices sont indépendants. Une lecture initiale de quelques minutes permet de repérer les thèmes, les questions directement accessibles, les questions de démonstration et les calculs plus longs. L’objectif n’est pas de prédire la difficulté de tout l’exercice, mais de construire un ordre de travail.

Une stratégie efficace consiste à marquer trois catégories : questions immédiatement accessibles, questions nécessitant une recherche courte et questions à reprendre plus tard. Cette classification évite de rester bloqué trop longtemps sur un point alors que des points indépendants restent disponibles ailleurs dans le sujet.

\coursefigure{/images/mathematiques/terminale-specialite-mathematiques/preparation-bac-specialite-mathematiques-1.svg}{Organisation indicative des quatre heures de l’épreuve.}{Une frise répartit lecture initiale, trois grandes phases de résolution et une réserve finale de vérification sur quatre heures.}

\section{2. Rédiger une réponse mathématique complète}
Une réponse attendue au bac contient, lorsque la question l’exige, quatre éléments : la propriété utilisée, ses hypothèses, le calcul ou le raisonnement, puis la conclusion. Pour un théorème des valeurs intermédiaires, on ne se contente pas d’écrire « par le TVI ». On indique la continuité, l’intervalle, les valeurs aux bornes et l’appartenance de la valeur cible à l’intervalle image.

De même, pour une étude de signe à partir d’une dérivée, on explique le signe du facteur déterminant avant d’en déduire les variations. Pour une probabilité, on définit l’événement ou la variable aléatoire avant d’utiliser une commande de calculatrice. La qualité rédactionnelle ne consiste donc pas à écrire beaucoup, mais à rendre chaque inférence vérifiable.

\section{3. Gérer un blocage}
Un blocage ponctuel ne doit pas interrompre toute la composition. On relit la question suivante : certains sujets donnent un résultat à admettre ou permettent de poursuivre avec une expression fournie. Il est alors légitime d’utiliser ce résultat pour continuer, tout en laissant une trace de la recherche précédente.

On peut aussi revenir à une forme plus simple du problème : écrire le domaine, calculer une dérivée, identifier une loi, dessiner une configuration ou tester une valeur. Ces micro-étapes peuvent débloquer le raisonnement et produisent souvent des éléments évaluables.

\section{4. Calculatrice : outil de calcul et de contrôle}
La calculatrice peut résoudre numériquement une équation, calculer une probabilité binomiale, produire une table de valeurs ou vérifier une valeur approchée. Elle ne justifie pas une existence, une unicité ou une propriété de variation.

Une bonne rédaction distingue donc le résultat théorique et le calcul numérique. Par exemple, on peut justifier qu’une équation possède une unique solution sur un intervalle par continuité et monotonie, puis utiliser la calculatrice pour donner une approximation de cette solution. L’approximation intervient après la preuve, pas à sa place.

\section{5. Analyse : chaîne de raisonnement type}
Un exercice d’analyse peut demander successivement domaine, dérivée, variations, limites, équation puis interprétation. Il est utile de conserver la chaîne logique :
\[
\text{domaine}\longrightarrow\text{dérivée}\longrightarrow\text{signe}\longrightarrow\text{variations}\longrightarrow\text{équation}.
\]
Chaque étape doit réutiliser les résultats précédents. Une erreur de domaine ou de signe au début peut contaminer toute la suite ; d’où l’intérêt des contrôles intermédiaires.

\section{6. Probabilités : modéliser avant de calculer}
Avant une loi binomiale, on justifie le schéma de Bernoulli : répétitions indépendantes, même probabilité de succès, nombre fixé d’essais. Pour une probabilité conditionnelle, on identifie clairement la condition. Pour une somme de variables, on distingue linéarité de l’espérance et additivité de la variance sous indépendance.

Une valeur calculée doit être interprétée dans le contexte. Écrire seulement une décimale ne répond pas à une question de risque, de fréquence attendue ou de seuil. La phrase de conclusion fait partie du raisonnement.

\section{7. Géométrie de l’espace}
Les exercices d’espace combinent souvent représentations paramétriques, équations cartésiennes, orthogonalité et distance. Une méthode robuste consiste à choisir des vecteurs directeurs ou normaux, vérifier les produits scalaires, puis traduire l’appartenance par substitution.

Pour une intersection entre une droite et un plan, on remplace les coordonnées paramétriques de la droite dans l’équation du plan. Pour une distance d’un point à un plan, on peut utiliser un projeté orthogonal ou une formule adaptée lorsque celle-ci est maîtrisée et justifiée.

\coursefigure{/images/mathematiques/terminale-specialite-mathematiques/preparation-bac-specialite-mathematiques-2.svg}{Carte de décision pour un exercice de spécialité.}{Un schéma relie analyse, probabilités, géométrie et algorithmique à leurs premiers réflexes de résolution et aux contrôles associés.}

\section{8. Calcul intégral et équations différentielles}
Dans un calcul intégral, on cherche une primitive, on contrôle les bornes et on distingue aire algébrique et aire géométrique. Pour une équation différentielle, on écrit d’abord la famille générale puis on utilise la condition initiale.

Ces deux thèmes se rencontrent dans des problèmes de modélisation. Une équation différentielle décrit un taux d’évolution, tandis qu’une intégrale cumule une grandeur. Le logarithme peut ensuite intervenir pour calculer un temps de franchissement.

\section{9. Algorithmique et traces de programme}
Lorsqu’un code est fourni, on identifie les variables, la condition de boucle, la mise à jour et la sortie. Une question peut demander ce que renvoie un programme, pourquoi il termine ou quelle propriété mathématique il traduit.

Pour une dichotomie, l’invariant est le maintien d’un changement de signe sur un intervalle où la fonction est continue. Pour une simulation, la fréquence obtenue est empirique. Le programme doit toujours être relié à l’objet mathématique qu’il représente.

\section{10. Vérification finale}
Les dernières minutes doivent être réservées au contrôle : domaines de définition, signes, unités, probabilités dans $[0,1]$, valeurs approchées cohérentes, bornes d’intégrales, réponses oubliées et conclusions. On vérifie aussi que chaque exercice traité comporte des réponses clairement repérables.

Une erreur de calcul isolée peut parfois être détectée par ordre de grandeur. Une probabilité supérieure à $1$, une durée négative dans un contexte impossible ou une aire négative sont des signaux immédiats. Le contrôle final doit donc être ciblé plutôt que purement esthétique.

\section{Méthode de travail avant l’épreuve}
L’entraînement efficace alterne trois échelles : exercices ciblés pour automatiser les techniques, exercices transversaux pour apprendre à choisir les outils et sujets complets en temps limité pour travailler l’endurance et la gestion du temps. Après chaque sujet, on classe les erreurs : connaissance manquante, méthode mal choisie, calcul, lecture de consigne ou gestion du temps.

Une fiche de révision utile n’est pas un catalogue de formules. Elle contient des déclencheurs de méthode : « logarithme : domaine d’abord », « TVI : continuité + encadrement », « loi binomiale : indépendance + même $p$ », « espace : vecteur normal ou directeur », « intégrale : primitive puis bornes ».

\section{Erreurs fréquentes}
\begin{itemize}
\item Rester trop longtemps sur une question alors que les exercices sont indépendants.
\item Donner une valeur calculatrice sans justification théorique lorsque la question demande une preuve.
\item Utiliser un théorème sans vérifier explicitement ses hypothèses.
\item Confondre résultat exact, valeur approchée et conjecture numérique.
\item Oublier une conclusion rédigée ou l’unité dans un problème contextualisé.
\item Ne pas exploiter un résultat fourni dans une question suivante après un blocage.
\end{itemize}

\section{À retenir}
\begin{aretenir}
L’épreuve valorise simultanément connaissances, choix des méthodes, qualité du raisonnement et communication. Les quatre exercices imposent une gestion active du temps. Une copie solide rend visibles les hypothèses, les calculs utiles, les contrôles et les conclusions.
\end{aretenir}
